% ============================================================================
% Lab Chapter for Part II — Practicals 1–8
% ============================================================================

\chapter{Part II --- Hands-On Practicals}
\label{ch:lab2}

\bytetip[fun]{Time to put your Windows skills to the test! These practicals cover everything from navigating the desktop to managing files like a pro.}

% === PRACTICAL 1 ===
\begin{labactivity}{1}{Exploring the Windows Desktop}{10 min}
\youwillneed{\faDesktop\ Computer \quad \faClock\ 10 min \quad \faFile\ Observation sheet}

\textbf{Objective:} Identify and interact with all desktop elements.\\[4pt]
\textbf{Steps:}
\begin{enumerate}[label=\protect\stepcircle{\arabic*}, leftmargin=2cm, itemsep=4pt]
  \item Identify the Start button, Taskbar, System Tray, and desktop icons.
  \item Click the Start button and explore the menu.
  \item Right-click on the desktop and choose ``Personalise.''
  \item Find the current time and date in the System Tray.
  \item Write down the names of all pinned apps on the Taskbar.
\end{enumerate}
\textbf{Expected Outcome:} Students identify all desktop elements correctly.
\welldone
\end{labactivity}

% === PRACTICAL 2 ===
\begin{labactivity}{2}{Changing Desktop Wallpaper}{10 min}
\youwillneed{\faDesktop\ Computer \quad \faClock\ 10 min}

\textbf{Objective:} Personalise the desktop appearance.\\[4pt]
\textbf{Steps:}
\begin{enumerate}[label=\protect\stepcircle{\arabic*}, leftmargin=2cm, itemsep=4pt]
  \item Right-click on an empty area of the desktop.
  \item Select ``Personalise'' from the context menu.
  \item Click ``Background'' and choose a new wallpaper from the list.
  \item Try selecting ``Solid colour'' and pick your favourite colour.
  \item Change it back to the original wallpaper when finished.
\end{enumerate}
\textbf{Expected Outcome:} Students successfully change and restore the wallpaper.
\welldone
\end{labactivity}

% === PRACTICAL 3 ===
\begin{labactivity}{3}{File and Folder Operations}{15 min}
\youwillneed{\faDesktop\ Computer \quad \faClock\ 15 min \quad \faFolder\ File Explorer}

\textbf{Objective:} Create, rename, copy, move, and delete files and folders.\\[4pt]
\textbf{Steps:}
\begin{enumerate}[label=\protect\stepcircle{\arabic*}, leftmargin=2cm, itemsep=4pt]
  \item Open \textbf{File Explorer} (\key{Windows}+\key{E}).
  \item Navigate to the \textbf{Documents} folder.
  \item Create a new folder: Right-click $\rightarrow$ New $\rightarrow$ Folder. Name it ``Practice.''
  \item Inside ``Practice,'' create two more folders: ``Text'' and ``Images.''
  \item Right-click inside ``Text'' $\rightarrow$ New $\rightarrow$ Text Document. Name it ``Hello.txt.''
  \item Right-click ``Hello.txt'' $\rightarrow$ Copy. Navigate to ``Images'' $\rightarrow$ Paste.
  \item Rename the copy to ``HelloCopy.txt'' (select it, press \key{F2}).
  \item Delete ``HelloCopy.txt'' (press \key{Delete}).
  \item Check the Recycle Bin --- is the file there?
\end{enumerate}
\textbf{Expected Outcome:} Students complete all file operations successfully.
\welldone
\end{labactivity}

% === PRACTICAL 4 ===
\begin{labactivity}{4}{Using Notepad}{10 min}
\youwillneed{\faDesktop\ Computer \quad \faClock\ 10 min \quad \faKeyboard\ Keyboard}

\textbf{Objective:} Create, edit, and save a text file using Notepad.\\[4pt]
\textbf{Steps:}
\begin{enumerate}[label=\protect\stepcircle{\arabic*}, leftmargin=2cm, itemsep=4pt]
  \item Open Notepad (Start $\rightarrow$ search ``Notepad'').
  \item Type: ``My name is [your name]. I am learning about Windows.''
  \item Press \key{Enter} and type today's date.
  \item Save the file: \key{Ctrl}+\key{S} $\rightarrow$ Save in Documents as ``MyNote.txt.''
  \item Close Notepad. Reopen ``MyNote.txt'' to verify your text is saved.
\end{enumerate}
\textbf{Expected Outcome:} Students save and retrieve a text file.
\welldone
\end{labactivity}

% === PRACTICAL 5 ===
\begin{labactivity}{5}{Drawing in Paint}{15 min}
\youwillneed{\faDesktop\ Computer \quad \faClock\ 15 min \quad \faPaintBrush\ Paint}

\textbf{Objective:} Use Paint's drawing tools to create a simple picture.\\[4pt]
\textbf{Steps:}
\begin{enumerate}[label=\protect\stepcircle{\arabic*}, leftmargin=2cm, itemsep=4pt]
  \item Open MS Paint (Start $\rightarrow$ search ``Paint'').
  \item Use the \textbf{Rectangle} and \textbf{Triangle} shapes to draw a house.
  \item Use the \textbf{Circle} tool to draw a sun.
  \item Add colour using the \textbf{Fill} tool (paint bucket).
  \item Use the \textbf{Text} tool to write your name at the bottom.
  \item Save as ``MyDrawing.png'' in Documents (\key{Ctrl}+\key{S}).
\end{enumerate}
\textbf{Expected Outcome:} A colourful drawing saved as a PNG file.
\welldone
\end{labactivity}

% === PRACTICAL 6 ===
\begin{labactivity}{6}{Using the Calculator}{10 min}
\youwillneed{\faDesktop\ Computer \quad \faClock\ 10 min \quad \faCalculator\ Calculator}

\textbf{Objective:} Perform calculations using the Windows Calculator.\\[4pt]
\textbf{Steps:}
\begin{enumerate}[label=\protect\stepcircle{\arabic*}, leftmargin=2cm, itemsep=4pt]
  \item Open Calculator (Start $\rightarrow$ search ``Calculator'').
  \item Calculate: 456 $\times$ 78. Write down the answer.
  \item Calculate: 9999 $\div$ 3. Write down the answer.
  \item Switch to \textbf{Scientific} mode (click the menu icon).
  \item Calculate $5^3$ (type 5, then use the $x^y$ button, type 3, press =).
  \item Switch to \textbf{Programmer} mode and explore binary numbers!
\end{enumerate}
\textbf{Expected Outcome:} Students use Standard and Scientific calculator modes.
\welldone
\end{labactivity}

% === PRACTICAL 7 ===
\begin{labactivity}{7}{Checking System Information}{10 min}
\youwillneed{\faDesktop\ Computer \quad \faClock\ 10 min \quad \faFile\ Worksheet}

\textbf{Objective:} Find and record basic system information.\\[4pt]
\textbf{Steps:}
\begin{enumerate}[label=\protect\stepcircle{\arabic*}, leftmargin=2cm, itemsep=4pt]
  \item Press \key{Windows}+\key{I} to open Settings.
  \item Go to \textbf{System} $\rightarrow$ \textbf{About}.
  \item Write down: Computer name, Processor, RAM amount, Windows version.
  \item Open \textbf{This PC} and check total storage on the C: drive.
  \item Record the total and free storage space.
\end{enumerate}
\textbf{Expected Outcome:} Students record their computer's specifications.
\welldone
\end{labactivity}

% === PRACTICAL 8 ===
\begin{labactivity}{8}{Taking Screenshots}{10 min}
\youwillneed{\faDesktop\ Computer \quad \faClock\ 10 min \quad \faCamera\ Snipping Tool}

\textbf{Objective:} Capture and save screenshots using different methods.\\[4pt]
\textbf{Steps:}
\begin{enumerate}[label=\protect\stepcircle{\arabic*}, leftmargin=2cm, itemsep=4pt]
  \item Press \key{PrtSc} to capture the full screen. Open Paint $\rightarrow$ \key{Ctrl}+\key{V} to paste.
  \item Save as ``FullScreen.png.''
  \item Press \key{Windows}+\key{Shift}+\key{S} to open Snipping Tool / Snip \& Sketch.
  \item Select a rectangular area of the screen.
  \item The snip is copied to clipboard --- paste it into Paint and save as ``MySnip.png.''
\end{enumerate}
\textbf{Expected Outcome:} Students capture and save two types of screenshots.
\welldone
\end{labactivity}

\vspace{12pt}
\begin{center}
  \qrcode[height=2cm]{https://example.com/module2}\\[6pt]
  {\small\sffamily\textcolor{NavyBlue}{Scan for extra Part II resources and activities!}}
\end{center}
